
\section{Mise en \oe{}uvre et conduite du projet en mode Agile}

Le projet a été conduit en s’inspirant des principes des méthodes Agile, sans appliquer strictement un cadre méthodologique comme Scrum. Cette approche nous a permis de conserver une certaine flexibilité tout en assurant un suivi régulier de l’avancement.

\subsection{Découpage en itérations}

Le développement a été organisé en plusieurs itérations correspondant à des sous-ensembles fonctionnels du projet. Une première itération a été consacrée à la mise en place des bases du langage MiniJaja (syntaxe, AST, mémoire), tandis que les itérations suivantes ont porté sur l’interprétation, la compilation vers JajaCode et l’intégration de l’interface.


\subsection{Suivi des tâches}

Les tâches ont été identifiées en début d’itération et ajustées au fil du projet en fonction des difficultés rencontrées. Certaines fonctionnalités ont nécessité plus de temps que prévu, en particulier lors de la phase d’intégration et de validation de la correction sémantique.


\subsection{Bilan de l’approche Agile}

L’approche Agile adoptée a permis de réagir rapidement aux problèmes techniques et d’adapter l’organisation du travail. Toutefois, le manque de formalisation de certains aspects (estimation précise des charges, planification détaillée) a parfois rendu le suivi plus difficile.

Globalement, cette méthode s’est révélée adaptée au contexte du projet et au cadre universitaire, en favorisant l’autonomie et la collaboration au sein du groupe.

